\documentclass[11pt,a4paper]{article}
\usepackage[margin=27mm]{geometry}
\usepackage{amsmath,amssymb,amsthm,hyperref}
\newtheorem{lemma}{Lemma}
\newtheorem{theorem}[lemma]{Theorem}
\newtheorem{corollary}[lemma]{Corollary}
\newcommand{\F}{\mathbb F}
\newcommand{\Sz}{\operatorname{Sz}}
\newcommand{\N}{\operatorname{N}}
\newcommand{\gen}[1]{\langle#1\rangle}
\title{Maximal intersections in Suzuki torus normalizers}
\author{Human Compression project}\date{}
\begin{document}\maketitle
\begin{abstract}
The subgroup definition of residual weak primitivity can be applied before
assuming flag transitivity. For cyclic groups and odd dihedral groups it
forces every parabolic to be maximal. In Suzuki groups, two successive
maximal intersections then reduce the possible torus normalizers to
compatible maximal subfields. Elementary size estimates exclude the
remaining pairs. A separate normalizer calculation handles two subfields
of the same size. These arguments prove the maximal-parabolic restrictions
in Lemmas 4.2--4.6 and the full split-normalizer case of Lemma 5.1 in the
fixed rank-three source.
\end{abstract}

\section{The subgroup condition in cyclic and dihedral groups}
For a family $(H_i)_{i\in I}$ in a finite group $H$, write
$H_J=\bigcap_{j\in J}H_j$ and $H_\varnothing=H$. We use the condition
\begin{equation}\label{rw}
 \text{for every }J\subsetneq I\text{ some }i\notin J
 \text{ makes }H_{J\cup\{i\}}\text{ maximal proper in }H_J.
\end{equation}
The previously accepted strict-intersection lemma says that adding
\emph{any} further type strictly decreases an intersection. Its proof is
short: $H_J\le H_i$ would imply
$H_{I\setminus\{i\}}=H_I$, contrary to \eqref{rw} at
$I\setminus\{i\}$. In particular no distinct parabolics contain one
another. The family induced inside any $H_J$ again satisfies \eqref{rw}.

\begin{lemma}\label{cyclic}
If $H$ is cyclic and the family satisfies \eqref{rw}, every $H_i$ is
maximal in $H$. Their indices are distinct primes.
\end{lemma}
\begin{proof}
Induct on $|I|$. For one type this is the definition. Choose $H_0$
maximal, of prime index $p$. For $i\ne0$, noncontainment in $H_0$
means that $|H_i|$ has the full $p$-part of $|H|$. Consequently
$|H_i\cap H_0|=|H_i|/p$. The induced family in $H_0$ satisfies
\eqref{rw}, so induction says that $H_i\cap H_0$ is maximal there.
Its index is
\[
 [H_0:H_i\cap H_0]=\frac{|H|/p}{|H_i|/p}=[H:H_i],
\]
which is prime. Equal indices would give equal subgroups in a cyclic
group, contrary to strict intersections.
\end{proof}

\begin{lemma}\label{dihedral}
Let $H=D_{2m}$ with $m$ odd. Every parabolic in a family satisfying
\eqref{rw} is maximal in $H$. At most one is cyclic. If the rank is
at least three, at least two are dihedral.
\end{lemma}
\begin{proof}
Again induct on the rank. Write $H=\gen r\rtimes\gen t$, with
$|r|=m$, $t^2=1$, and $r^t=r^{-1}$. If a chosen maximal parabolic
is $A=\gen r$, every other parabolic is dihedral by noncontainment.
Its intersection with $A$ is its rotation subgroup, which is maximal
in $A$ by Lemma~\ref{cyclic} applied to the induced family. The
parabolic therefore has prime index in $H$.

Otherwise choose $H_0=\gen{r^p,t}$ for a prime $p\mid m$, after
changing the reflection $t$. Every $H_0\cap H_i$ is maximal in
$H_0$ by induction (or Lemma~\ref{cyclic} when $H_0=C_2$).
If $H_i=C_d$ is cyclic, noncontainment in $H_0$ gives the full
$p$-part in $d$, and its intersection has order $d/p$.
The only maximal cyclic subgroup of $H_0$ is its rotation subgroup,
so $d/p=m/p$ and $H_i=A$.

If $H_i=\gen{r^{m/d},r^k t}$ is dihedral, first suppose $d$
has the full $p$-part of $m$. Then $p\nmid m/d$. The two reflection
congruences have a simultaneous solution, so $H_0\cap H_i$ is
dihedral of order $2d/p$. Its index in $H_0$ is $m/d$; since it is
maximal, that index is prime. Thus $H_i$ is maximal in $H$.
If $d$ has a smaller $p$-part, then $d\mid m/p$. Noncontainment
forces its reflection coset to be disjoint from $H_0$, so
$H_0\cap H_i=C_d$. Maximality in $H_0$ forces $d=m/p$, and
again $H_i$ is maximal, now of index $p$.
This covers every subgroup. A maximal cyclic subgroup must be $A$,
so there is at most one; the last assertion follows immediately.
\end{proof}

This proves directly the dihedral consequence cited from reference~[20]
in the source's Lemmas 4.3 and 4.4. It does not claim a new classification
of all dihedral geometries or import their unexplained incidence properties.

\section{Structural inputs and a subfield uniqueness lemma}
Put $q=2^n$, with $n\ge3$ odd. We retain the classical Suzuki subgroup,
ovoid and matrix description specified in the accepted L115 notes.
Every proper subgroup either lies in a point stabilizer, a split torus
normalizer, a nonsplit torus normalizer, or is itself a simple subfield
group. Write
\[
 B_q=U_q\rtimes\F_q^*,\quad D_q=C_{q-1}\rtimes C_2,
 \quad N_\epsilon(q)=C_{a_\epsilon(q)}\rtimes C_4,
 \quad a_\epsilon(q)=q+\epsilon\sqrt{2q}+1.
\]
The three odd torus orders are pairwise coprime, and a nonidentity odd
element lies in a unique torus. A proper simple subfield size satisfies
$q\ge s^3$. The accepted elementary estimates imply, also directly,
\begin{equation}\label{bounds}
 q/2\le a_\epsilon(q)<2q \quad(q\ge8).
\end{equation}
Indeed $q+1-\sqrt{2q}\ge q/2$ and
$q+1+\sqrt{2q}<2q$ in this range.
The uniqueness of an embedded $D_{2(s-1)}$ or
$C_{a_\delta(s)}\rtimes C_4$ inside a fixed-size subfield group
was proved in the accepted normalizer-fibre note.

\begin{lemma}\label{affine}
For $s=2^d\ge8$, $d\mid n$, every subgroup conjugate to the standard
$Z_s\rtimes\F_s^*$ is contained in exactly one subgroup isomorphic
to $\Sz(s)$. Here $Z_s$ is the centre of a Sylow two-subgroup of
the standard subfield group.
\end{lemma}
\begin{proof}
We prove that the standard $H=Z_s\rtimes\F_s^*$ is self-normalizing
in $G=\Sz(q)$. Its characteristic Sylow two-subgroup $Z_s$ determines
its unique fixed ovoid point. A normalizer therefore lies in $B_q$.
Use the accepted coordinates $u(a,b)$ for its Sylow subgroup and $h(v)$
for the diagonal torus. Conjugation by $h(v)$ multiplies the central
coordinate by $v^{\theta+2}$, where $\theta=2^{(n+1)/2}$.
Preserving the central line $\F_s$ forces
$v^{\theta+2}\in\F_s^*$. Since $\theta+2$ is coprime to $q-1$,
this is equivalent to $v\in\F_s^*$. Remove this diagonal factor.

For $u(a,b)$, preservation of $H$ forces every commutator with
$h(k)$, $k\in\F_s^*$, to have first coordinate zero. That coordinate
is a nonzero scalar multiple of $(k^\theta-1)a$. Taking $k\ne1$
therefore gives $a=0$. The remaining commutator coordinate is a
nonzero scalar multiple of $(k^{\theta+2}-1)b$ and must lie in
$\F_s$. Again take $k\ne1$; its coefficient is a nonzero element
of $\F_s$, so $b\in\F_s$. Thus $\N_G(H)=H$.

The standard subfield group $S=\Sz(s)$ is self-normalizing, as proved
in the accepted note. Its copies of the indicated affine subgroup form
one conjugacy class: the Sylow centres are conjugate, and the odd
complements in their point stabilizers are conjugate by the elementary
affine-complement argument. The same incidence count used in that note
therefore gives
\[
 \#\{S^g:H\le S^g\}=\frac{|\N_G(H)|}{|\N_S(H)|}=1.
\]
Conjugation gives the assertion for every conjugate of $H$.
\end{proof}

We also need an elementary overgroup observation valid for every $s$.
In $S=\Sz(s)$, the proper overgroups of its full split torus
$A_s=C_{s-1}$ are exactly
\begin{equation}\label{fullsplit}
 A_s,\quad D_s,\quad Z_P A_s,\quad Z_Q A_s,\quad B_P,\quad B_Q,
\end{equation}
where $P,Q$ are the two fixed ovoid points. The classical list excludes
nonsplit normalizers by coprimality and proper simple subfields by their
smaller torus orders. Inside either point stabilizer an overgroup is
$V A_s$. Full scalar multiplication acts irreducibly on both $Z(U)$
and $U/Z(U)$ over $\F_2$. Thus $V=1,Z(U)$ or $U$: if its image
on $U/Z(U)$ is nonzero, then $VZ(U)=U$, and the accepted identity
$Z(U)=\Phi(U)$ gives $V=U$. This proves \eqref{fullsplit}.

\section{Excluding a full nonsplit normalizer}
\begin{theorem}\label{nonsplit}
A family satisfying \eqref{rw} in $\Sz(q)$ cannot have a full
nonsplit torus normalizer as a parabolic when its rank is at least three.
\end{theorem}
\begin{proof}
The prime-degree case is already accepted, so assume $n$ composite.
Both nonsplit torus orders are then composite by the accepted subfield
arithmetic. Suppose $G_0=N=C_m\rtimes C_4$, with $m=a_\epsilon(q)$.
Choose $J=G_{01}$ maximal in $N$. Its possibilities are
$D_{2m}$ and $C_{m/p}\rtimes C_4$, $p\mid m$ prime.
No proper subgroup outside $N$ contains $D_{2m}$: the classical list
and the torus partition exclude the other solvable types, and a proper
simple subfield has every torus order less than $2s<q/2\le m$.
The first possibility is therefore excluded.

Set $d=m/p>1$ and $J=C_d\rtimes C_4$. The proper overgroup $G_1$
is not contained in $N$ by strict intersections. The classical list
forces $G_1=\Sz(s)$. Inside it the odd subgroup $C_d$ lies in a
nonsplit torus of order $a_\delta(s)$, which divides $m$ and is
smaller than $m$. Since it is a multiple of $m/p$, it equals $d$.
Thus $J$ is its full normalizer in $G_1$. Moreover $G_1$ is a maximal
subfield group: an intermediate field would supply a strictly intermediate
compatible torus order between $d$ and $m$, impossible when $m/d=p$
is prime. Write $q=s^a$ with $a$ an odd prime.

Choose another type $2$ so that $K=G_{012}$ is maximal in $J$.
If $K=D_{2d}$, then $G_{02}$ contains it. A subgroup of $N$ containing
$K$ has odd kernel of order $d$ or $m$. In the first case it is contained
in the unique $C_d\rtimes C_4$ containing $K$, namely $J$.
In the second case it contains $D_{2m}$, which cannot be in the proper
subgroup $G_2\nleq N$. Both contradict strict intersections.

The other possibility is $K=C_{m/(pr)}\rtimes C_4$ for a prime
$r\mid d$. Since $G_{02}$ contains this fixed $C_4$, its odd kernel
has order divisible by $m/(pr)$ and dividing $m$. Noncontainment in
$J$ requires the full $p$-part. If $r=p$, this forces the whole $N$,
which is impossible. If $r\ne p$, it forces
$G_{02}=C_{m/r}\rtimes C_4$. The same argument as for $G_1$ gives
$G_2=\Sz(t)$, with $m/r=a_\eta(t)$ and $q=t^b$, $b$ prime.

If $s=t$, the two torus orders are distinct and hence are the opposite
ones. Their product is $s^2+1$, whereas
\[
 m=\frac{(m/p)(m/r)}{m/(pr)}\le s^2+1.
\]
This contradicts $m\ge q/2\ge s^3/2>s^2+1$.
If $s\ne t$, the distinct maximal field indices $a,b$ give
$q=2^{hab}$, $s=2^{hb}$ and $t=2^{ha}$ for a positive integer $h$.
Distinct odd primes satisfy $ab-a-b\ge7$, so $q/(st)\ge128$.
But $p$ divides $m/r=a_\eta(t)$, while \eqref{bounds} gives
\[
 p=\frac{m}{a_\delta(s)}>\frac{q}{4s}\ge32t>a_\eta(t).
\]
This final contradiction proves the theorem.
\end{proof}

\section{The full split normalizer}
\begin{theorem}\label{split}
Suppose $q-1$ is composite and a family of rank at least three satisfying
\eqref{rw} has $G_0=D_q$. At most one $G_{0i}$ is dihedral.
\end{theorem}
\begin{proof}
Apply Lemma~\ref{dihedral} to the family induced in $G_0$.
Every $G_{0i}$ is maximal in $D_q$. Suppose two are dihedral:
$G_{01}=D_{2m/p}$ and $G_{02}=D_{2m/r}$, where $m=q-1$ and
$p,r$ are primes. Their rotation subgroups are nontrivial.
A proper overgroup of such a split dihedral subgroup cannot lie in a
point stabilizer: its involution has trivial image in the abelian odd
quotient, so cannot invert the nontrivial odd image of a rotation.
Coprimality excludes nonsplit normalizers, and the torus partition
identifies the only possible split normalizer as $D_q$.
Consequently $G_1=\Sz(s)$ and $G_2=\Sz(t)$.
Their split torus orders are multiples of $m/p$ and $m/r$, respectively,
and proper divisors of $m$. Hence
\[
 s-1=m/p,\qquad t-1=m/r.
\]
As in the preceding proof, the prime indices imply that both are
maximal subfield groups.

If $p\ne r$, then $s\ne t$. Write their distinct field indices as
$a,b$. The same calculation gives $q/(st)\ge128$.
The prime $p$ divides $m/r=t-1$, but
\[
 p=\frac{q-1}{s-1}>\frac{q}{2s}\ge64t>t-1,
\]
a contradiction.

If $p=r$, the two distinct dihedral subgroups have the same rotation
subgroup $A_s=C_{s-1}$ and disjoint reflection cosets. Thus
$G_{012}=A_s$, and $s=t$. Strictness requires
$L=G_1\cap G_2>A_s$. Apply \eqref{fullsplit} inside $G_1$.
If $L=D_s$, it equals the normalizer of $A_s$ in $G_1$, namely
$G_{01}$. Then $G_{012}=G_{01}$, a contradiction. Each other
possibility contains one of the groups $Z_P A_s$ or $Z_Q A_s$.
By Lemma~\ref{affine}, this subgroup is contained in a unique copy
of $\Sz(s)$, so $G_1=G_2$, again impossible. This treats the
same-size case separately; distinct field indices were not assumed
merely from distinct parabolics.
\end{proof}

\begin{corollary}\label{rank}
A family satisfying \eqref{rw} with a full split normalizer as a
parabolic has rank at most three. In rank at least four, every
parabolic maximal in $\Sz(q)$ must be a simple subfield group.
\end{corollary}
\begin{proof}
If $q-1$ is prime, the earlier accepted rank bound applies.
If it is composite and the rank is at least four, the residue in
$D_q$ has rank at least three. Lemma~\ref{dihedral} gives at least
two dihedral parabolics in it, contrary to Theorem~\ref{split}.
For the last assertion, the accepted point-stabilizer exclusion and
Theorem~\ref{nonsplit} exclude the other solvable maximal subgroups.
The classical maximal-subgroup list leaves only simple subfields.
\end{proof}

\begin{corollary}\label{shape}
Suppose $n$ is composite and a rank-three family satisfying \eqref{rw}
has $G_0=D_q$. After interchanging types 1 and 2, its pattern is
\[
\begin{array}{c|ccc}
 &G_0&G_1&G_2\\\hline
 &D_{2(q-1)}&E_q\rtimes C_{q-1}&\Sz(s)
\end{array}
\]
\[
 G_{01}=C_{q-1},\quad G_{02}=D_{2(s-1)},\quad
 G_{12}=E_s\rtimes C_{s-1},\quad G_{012}=C_{s-1},
\]
where $(q-1)/(s-1)$ is prime. The degree $n$ is a power of one
odd prime. This is a necessary pattern, not an existence assertion.
\end{corollary}
\begin{proof}
The two proper intersections in $D_q$ are maximal by
Lemma~\ref{dihedral}. At most one is dihedral by Theorem~\ref{split},
and there cannot be two cyclic ones. Thus one is $C_{q-1}$ and the
other is $D_{2(q-1)/p}$. The proof of Theorem~\ref{split} gives
$G_2=\Sz(s)$ with $s-1=(q-1)/p$.
The full-torus overgroups \eqref{fullsplit}, now in $\Sz(q)$, and
strictness leave $G_1=Z_P C_{q-1}$ or $Z_Q C_{q-1}$; the full
point stabilizers have already been excluded. Their intersection with
$G_2$ properly contains $C_{s-1}$. Applying \eqref{fullsplit} inside
$G_2$, it cannot be dihedral, since a split dihedral group cannot lie
in a point stabilizer. It cannot be the full $B_s$, whose Sylow
two-subgroup has exponent four, whereas every two-subgroup of $G_1$
is elementary abelian. The remaining option is $E_s\rtimes C_{s-1}$.

For completeness, the prime quotient forces the degree assertion.
Write $q=s^a$. If $a$ were composite, the geometric sum
$(s^a-1)/(s-1)$ would factor into two integers greater than one;
thus $a$ is prime. Write $s=2^d$. If a prime $b\ne a$ divided $d$,
put $x=2^{d/b}$. The smaller integer
$M=1+x+\cdots+x^{a-1}>1$ divides
$1+x^b+\cdots+x^{b(a-1)}$: modulo $M$, $x^a=1$, and multiplication
by $b$ permutes the residues modulo $a$. This would again make
$(q-1)/(s-1)$ composite. Therefore $d$ is a power of $a$, and so
is $n=ad$.
\end{proof}

\section{Attribution and scope}
The source is D.~Leemans, \emph{The rank 3 geometries of the simple
Suzuki groups $\Sz(q)$}, Note di Matematica 19(1) (1999), 43--63.
Theorem~\ref{nonsplit} proves its Lemma~4.2 for every permitted field
degree. Lemma~\ref{dihedral}, Theorem~\ref{split} and
Corollary~\ref{rank} prove Lemmas~4.3--4.6; Corollary~\ref{shape}
proves the precise necessary assertion in Lemma~5.1.
The cyclic and dihedral induction replaces the explicitly cited
consequence of the earlier dihedral classification. The same-size
subfield argument is supplied explicitly rather than inferring distinct
field degrees from distinct subgroups. These are elementary arguments;
no priority claim is made for the methods.

The classical subgroup classification, the faithful Suzuki coordinates,
ovoid action and subfield embeddings remain borrowed structural inputs.
The accepted L115 notes supply the Frattini identity, scalar actions,
torus arithmetic and normalizer-incidence counting. The accepted L114
prime-degree proof supplies strict intersections and the point-stabilizer
and prime-degree exclusions. No finite computation or Lean certificate
is asserted in this component. It does not repair the refuted complete
candidate list in Theorem~5.5, prove the remaining thinness assertion,
or confer whole-original or new-conjecture credit.
\end{document}
